% P10-4a, 14.09.2026: Mechanismus, Quellenreichweite und v4-Nachweis angeglichen.
% M17, 12.09.2026: Gesamtfluss und Ebenen lesbarer getrennt; Grafik, Tabellen und Governance-Reichweiten erhalten.
% ============================================================
% §5.2 Logische Gesamtarchitektur — v06 (09.09.)
% v06: Erst-Seed ohne Einzelautorisierung von den nachgelagerten
% Strukturierungs- und Fortschreibungspfaden getrennt; Text, Caption
% und Eingangstabelle an 5.6/5.9/7.4 angeglichen. Grafik separat ändern.
% v05 (Gesamtdurchgang P0/P1, todo-v02 — Mini-Touches, keine Logik-
% Änderung): ① Zugangsformen auf Konvergenz-Endform (P0-2, wie
% 4.9-v07/5.8) ② damalige Bootstrap-Formulierung in v06 durch die
% belegte Trennung von Erst-Seed und nachgelagerten Gates ersetzt
% ③ „Abkürzung existiert nicht" → Architekturvertrag
% des vorgesehenen Pfads (P1; 5.6-Grenze: Naht beweist Gate-Herkunft
% nicht) ④ Tabellen-Zeile Meeting: „SELEKTIVER menschlicher
% Adjudikation bei definiertem Prüfbedarf" (P1; kein Gate pro Claim).
% MERKER an den AUTOR: In G1 den Erst-Seed als eigenen Weg zur
% Speichernaht ohne Einzelautorisierung zeigen; nachgelagerte Themen-
% und Arbeitspaket-Vorschläge durchlaufen die regulären Gates.
% v04 (Abschluss-Review, beide Mikro-Punkte übernommen): ① „Alles,
% was nach einer Zustandsänderung geschieht" → „Die nachgelagerten
% Fortschreibungs- und Projektionspfade arbeiten auf diesem geltenden
% Zustand" (konsequente Anwendung der Trigger-Logik-Klärung: kein
% „Save löst Folgeprozess aus") ② „Schichtung von Architekturebenen"
% → „Gliederung in Steuerungs-, Ausführungs- und Zustandsebene".
% P6-„können"-Variante des Kollegen bewusst NICHT übernommen (er
% selbst: „notwendig ist diese Änderung aber nicht") — Indikativ hält
% die Wort-Parität zu Kap. 4.9. Kollegen-Urteil: „inhaltlich
% geschlossen", FROZEN-Empfehlung — Freeze-ENTSCHEID = Autor.
% v03 (Kollegen-Review, kritisch geprüft — P1/2/3/5/7 übernommen,
% P4 abgeschwächt übernommen, P6 NUR TEILWEISE): ① Bootstrap in
% Lebenszyklus-Absatz + Meeting-Tabellenzeile (P1 — echtes Loch,
% repo-wahr: Erstbestand ohne Einordnung) ② Klärungsantwort als
% „interner Wiedereintritt" ausgewiesen (P2 — Karte §0: kein vierter
% Außen-Eingang; Tabellen-Caption/Spaltenkopf angepasst) ③ lokale
% Ableitung vs. gated Außenkopplung getrennt (P3 — GithubDocPublish-
% verifiziert) ④ Schleifen-Satz minimal entschärft ⑤ „Architektur-
% ebenen" statt „Verantwortungsebenen" + „getrennt geführt" (P5;
% bewusste Abweichung vom Blaupausen-Wort — Nachzug klein) ⑥ Steward-
% Grammatik + „Schreibweg" (P6) — ABER „bedienen dieselben
% Entscheidungspunkte" behalten (Kap.-4.9-Wortlaut + K14; Kollegen-
% Abschwächung „stärker als notwendig" abgelehnt, weil repo-wahr)
% ⑦ Einstieg geglättet (P7).
% v02 (nach Grafik-FINAL gesamtarchitektur-reduziert-final.drawio,
% Kollegen-Abnahme „freeze-ready"): ① figure-TODO auf finalen
% Dateinamen ② Viererketten-Wortlaut an Grafik angeglichen
% („pfadspezifische Prüfung", „fachliche Autorisierung", „autorisierte
% Mutation") ③ Brücken-Satz zur zusammengefassten Eingangsaufbereitung
% (funktionale Verantwortung ≠ gemeinsame Implementierung) ④ Kommentar
% Zugangs-Bindung aktualisiert (Steward-Steuerungspfeil jetzt in der
% Grafik). Der Core-Satz „ausschließlich Zustand … keine Verarbeitung"
% ist seit dem Grafik-Kürzen (Kollegen-Punkt 8) TRAGEND im Text —
% nicht streichen!
% Grundlage: Schreibblaupause-Kapitel-5-6 §4 (Zeile 5.2: ~1,5 S.;
% Pflichtinhalte: Gesamtgrafik + zwei Ansichten + Zustandsarten-
% Trennung + Eingangs-Tabelle + Zugangsformen/Ein-Kette + G2) ·
% architekturkarte-5-2.md v04 (FROZEN; §1b Lesehilfe = Rohmasse,
% Kantenmatrix K1–K16 = Beleg-Grundlage) · Grafik-Abnahme 02.09.
% (reduzierte Version + Detailversion konsistent; H an Doku-Sichten
% per Autor-Beschluss 02.09.).
%
% PRÄZISIONS-PUNKTE (heutige Klärungsrunde, alle code-verifiziert):
%   · Fortschreibung wird vom GELTENDEN ZUSTAND gespeist, nicht vom
%     Speicher-Ereignis (keine Trigger-Logik in der Naht; K9/K10).
%   · Schleifen-Ergebnisse münden als typisierte Vorschläge in die
%     Governance — KEINE erneute Fünf-Kategorien-Einordnung
%     (E47-02: Gate-Auflösung → DecisionApply direkt).
%   · Beide Projektionsarten menschlich freigegeben, aber EINE
%     Gate-Instanz (GithubDocPublish über dasselbe Forward-Gate).
%   · Zugangs-Bindung nicht exklusiv: nur die dialogische Eingabe
%     ist steward-exklusiv; beide Zugänge starten dieselben Abläufe.
% EINMAL-PRINZIP: Viererkette hier nur BENANNT (Vertiefung 5.6);
% Einordnungs-Kategorien nur benannt (5.5); Evidenzgewinnung (5.3);
% Projektions-Rundweg (5.7); Durabilität/Steward (5.8). KEINE
% Mechanismus-Erklärung in 5.2 — die Grafik-Führung zeigt nur, WAS
% wo liegt und WIE es zusammenhängt.
% Verweis-Merker: Refs auf 5.3–5.8 stehen als Klartext-Nummern und
% werden nach Anlage der Abschnitte auf \ref umgestellt (Label-Regel
% wie 3.2). Abbildung: Platzhalter-figure — Autor bindet den Export
% der reduzierten drawio-Version ein (Detailversion → Anhang).
% Budget ~1,5 S. + zwei kompakte Tabellen.
% ============================================================

\section{Logische Gesamtarchitektur}
\label{sec:logische-gesamtarchitektur}

Abbildung~\ref{fig:logische-gesamtarchitektur} zeigt den Zusammenhang
der Architekturbausteine; eine detailliertere Fassung steht in
Anhang~\ref{anh:g1-detail}. Zwei Perspektiven erschließen die
Darstellung: Der funktionale Lebenszyklus beschreibt, wie
Projektinformationen verarbeitet werden. Die Architekturebenen
unterscheiden Steuerung, Ausführung und fachlichen Zustand.

\begin{figure}[htbp]
  \centering
  % TODO (Autor): finalen Export einbinden —
  % chap5/gesamtarchitektur-reduziert-final.drawio.pdf, z. B.
  \includegraphics[width=\textwidth]{figures/05/GesamtArchitektur-Reduziert-v02.pdf}
  \caption[Logische Gesamtarchitektur des Artefakts]{Logische Gesamtarchitektur des Artefakts
    (Überblicksdarstellung; H = menschlicher Entscheidungspunkt).
    Der Erst-Seed schreibt den ersten fachlichen Bestand über die
    gemeinsame Speicheroperation (Speichernaht) ohne Einzelautorisierung
    (Abschnitt~\ref{sec:endtoend-pfade}).
    Detailfassung im Anhang.}
  \label{fig:logische-gesamtarchitektur}
\end{figure}
% P10-6: Gesamtgrafik vor ihrer ausführlichen Leseanleitung ausgeben.
\clearpage

\emph{Erste Ansicht: der funktionale Lebenszyklus.} Der Hauptfluss
verläuft von oben nach unten. Drei externe Eingänge liefern neue
Projektinformationen: Transkripte, dialogische Eingaben und
GitHub-Rückmeldungen (Tabelle~\ref{tab:eingaenge}). Sie unterscheiden
sich zunächst in ihrer Aufbereitung. Beim Transkript führen
Segmentierung, Claim-Extraktion und fachliche Ableitung zu
strukturierten Inhalten. Der dialogische Eingang übernimmt die im
Steward-Gespräch formulierten Angaben; beim GitHub-Eingang bereiten
Agenten die eingesammelten Rückmeldungen fachlich auf. Die jeweiligen
Schritte halten die verfügbaren Herkunftsangaben fest. Ihr gemeinsames
Ergebnis ist ein \emph{Delta}: ein strukturiertes Eingangspaket mit
neu zu prüfenden Inhalten und Quellenbezügen, noch kein autorisierter
Änderungsplan.

Bei vorhandenem Core münden diese Pakete in die gemeinsame
Einordnung gegen den Bestand, mit fachabhängigen Zweigen für
Anforderungen und Architekturinformationen. Erst dort entstehen
konkrete Operationspläne: Sie benennen die vorgeschlagene
Zustandsänderung und gegebenenfalls ihr bestehendes Ziel. Der
Governance-Bereich verbindet die pfadspezifische Prüfung dieser
Vorschläge mit menschlicher Autorisierung, Anwendung und Speicherung.
Die gemeinsame Speicheroperation, in der Abbildung als
\emph{Speichernaht} bezeichnet, prüft strukturelle Invarianten.
Über diesen Weg werden freigegebene Änderungen Teil des Core, des
maßgeblichen Projektzustands über einzelne Läufe hinweg. Die
Begründung dieser Kontrollen steht in
Abschnitt~\ref{sec:projektzustand-kontrollierte-mutation}, ihre
Ausgestaltung in Abschnitt~\ref{sec:governance-mutation}.

Der initiale Meetingpfad stellt dagegen den ersten Bestand her.
Aus der verwendbaren Evidenzsicht nach selektiver Adjudikation
entstehen modellgestützt fachliche Artefakte. Ein deterministischer
\emph{Erst-Seed} übernimmt daraus Anforderungs- und Architekturelemente
über die invariantengeschützte Speichernaht in den erstmals
angelegten Core. Für diese Elemente fehlt eine gesonderte
menschliche Einzelautorisierung. Die anschließende Themenstrukturierung
und Arbeitspaketbildung durchlaufen die vorgesehenen menschlichen
Gates. Das hergeleitete Autorisierungsziel ist bei der Initialisierung
somit nur teilweise realisiert. Abschnitt~\ref{sec:endtoend-pfade}
erklärt die Reihenfolge der Schreibvorgänge und Entscheidungen.

Der Core enthält Elemente, Relationen, Statusangaben, Historie
und Herkunftsinformationen. Die Verarbeitung leisten die
umgebenden Komponenten. Die beiden Fortschreibungspfade unterhalb
des Core-Bandes bearbeiten offene Entscheidungen und bereiten die
Angleichung abhängiger Arbeit vor. Sie nutzen den gespeicherten
Bestand einschließlich offener Vorgänge aus früheren Läufen.
Der Bezug zum Bestand steht hier bereits fest; eine erneute
Einordnung neuer Informationen ist deshalb nicht Teil ihrer Aufgabe.
Die Klärungsantwort bildet einen solchen begrenzten internen
Wiedereintritt. Die resultierenden Änderungsvorschläge gelangen
wiederum in den Governance-Bereich und verändern den fachlichen
Bestand nicht unmittelbar.

Aus dem Core entstehen schließlich operative Sichten:
GitHub-Issues mit eindeutiger Bindung an Core-Elemente sowie
maschinell gepflegte Dokument- und Architektursichten. Ihre lokale
Ableitung ist von der Veröffentlichung zu unterscheiden. Die
Überführung nach außen wird für beide Projektionsarten über
dieselbe menschliche Autorisierungsinstanz freigegeben. Nur der
GitHub-Weg besitzt einen Rückkanal: Eingesammelte Außenänderungen
gelangen erneut über einen kontrollierten Eingang in die
Verarbeitung (Abschnitt~\ref{sec:projektionen-kontrollierte-aussenkopplung}).

\begin{table}[htbp]
  \centering
  \small
  \begin{tabular}{p{0.24\textwidth}p{0.40\textwidth}p{0.26\textwidth}}
    \hline
    \textbf{Pfad} & \textbf{Aufbereitung} & \textbf{Eintritt} \\
    \hline
    Meeting-Transkript & Evidenzgewinnung: quellgebundene Belegeinheiten
      mit selektiver menschlicher Adjudikation bei definiertem
      Prüfbedarf und anschließender fachlicher Ableitung
      (Abschnitte~\ref{sec:evidenzgebundene-verarbeitung}
      und~\ref{sec:saeule-semantische-transformation}) & Bootstrap: Erst-Seed und
      nachgelagerte Strukturierung; Betrieb: Einordnung gegen den
      Bestand \\
    GitHub-Rückmeldung & modellgestützte Aufbereitung der eingesammelten
      Außenänderungen, mit Herkunftsangabe & Einordnung gegen den
      Bestand \\
    dialogische Eingabe & kontrollierter Entwurf über die agentische
      Bedienschicht, mit Herkunftsangabe & Einordnung gegen den
      Bestand \\
    interner Wiedereintritt: Klärungsantwort & kontrollierte Auflösung
      am menschlichen Entscheidungspunkt; eng begrenztes
      Auflösungsvokabular & direkt in die Governance \\
    \hline
  \end{tabular}
  \caption{Kontrollierte Eingangspfade und interner Wiedereintritt.}
  \label{tab:eingaenge}
\end{table}
% [Zeile 1 = Kollegen-P1-Konsequenz (Bootstrap/Betrieb). Zeile 3
%  (dialogische Eingabe) = K6, GEHÄRTET 02.09.: Lauf runs/fullworkflow/
%  20260817_123342_f23f06 — Diktat im AF-Delta → Tor 1, reviewer
%  „author via steward-chat", AF-1 apply → REQ-42→v2 + PBI-028→v2.
%  Zeile 4 = heutige Klärungsrunde: die Antwort am Decision-Gate ist
%  ein eng begrenzter Wiedereintritt (KEEP/ADOPT_NEW/REFINE/
%  NO_TRUTH_NEEDED, nur Ziel-Element) — mehr implizierte Wahrheit
%  nimmt den normalen Eingangsweg. Karte §1b/⑤.]

\emph{Zweite Ansicht: Architekturebenen.} Die Steuerungsebene
umfasst die Zugangsformen. Die Ausführungsebene verbindet die
Verarbeitungsschritte zu einem planmäßig unterbrechbaren Ablauf
mit speicherbarem Ausführungszustand. Der Core führt davon getrennt
den fachlichen Projektzustand. Der Ausführungszustand beantwortet,
wo ein Lauf steht und welche Entscheidung aussteht; der Core hält
fest, was im Projekt gilt. Diese in
Abschnitt~\ref{sec:durabilitaet-kontinuierliche-nutzung} begründete
Trennung erläutert Abschnitt~\ref{sec:durable-ausfuehrung-agentische-bedienung}
mit ihren Folgen für den Betrieb.

\emph{Zugangsformen.} Von den fachlichen Eingängen sind die
Arten der Bedienung zu unterscheiden. Der Direktzugang startet
vorgegebene Abläufe; der \emph{Steward} interpretiert Anliegen im
Dialog, wählt passende Systemfähigkeiten und vermittelt anstehende
Entscheidungen. Beide nutzen für ihre fachlichen Wirkungen die
vorgesehenen Prüf-, Entscheidungs- und Speicheroperationen. Die
dialogische Eingabe ist an den Steward gebunden; dieser kann auch
andere vorhandene Abläufe bedienen und besitzt keinen eigenen
Core-Schreibweg. Seine Rolle vertieft
Abschnitt~\ref{sec:durable-ausfuehrung-agentische-bedienung}.

\emph{Verantwortungsverteilung.} Tabelle~\ref{tab:verantwortung}
ordnet die Aufgaben den primär zuständigen Komponenten oder dem
Menschen zu. Die folgenden Säulenabschnitte zeigen, wie diese
Verantwortlichkeiten im jeweiligen Verarbeitungspfad zusammenwirken.

\begin{table}[htbp]
  \centering
  \small
  \begin{tabular}{p{0.56\textwidth}p{0.32\textwidth}}
    \hline
    \textbf{Aufgabenart} & \textbf{primäre Verantwortung} \\
    \hline
    Bedeutung interpretieren, einordnen, formulieren, semantisch
      prüfen oder reparieren & agentisch (modellgestützt) \\
    formal ausdrückbare Prüfung, Routing, Vertrags- und
      Integritätsregeln, autorisierte Anwendung & deterministisch \\
    fachlich folgenreiche Entscheidungen und Autorisierung & Mensch \\
    Reihenfolge, Wiederholung, Verzweigung, Übergaben, Pausen &
      Workflow \\
    \hline
  \end{tabular}
  \caption{Verantwortungsverteilung in der Gesamtarchitektur.}
  \label{tab:verantwortung}
\end{table}
% [G2 aus der Karte §5 (Blaupause: „G2 in 5.2 oder 5.4" — hier, weil
%  sie die Lesebrille für ALLE Säulen-Zooms ist; „dosierte Agentik").]

Diese Verteilung erhält agentischen Entscheidungsspielraum bei der
Interpretation und Werkzeugwahl. Der Workflow legt dagegen fest,
wann ein Ergebnis geprüft, zur Entscheidung vorgelegt und angewendet
wird. Die in Kapitel~\ref{chap:explorative-problemklaerung}
beobachteten Schwierigkeiten bei Review, Toolausführung und
Zustandsfortschreibung motivieren diese Bindung an überprüfbare
Zwischenstände. Sie ist eine begründete Gestaltung für den
untersuchten Einsatz, keine nachgewiesene allgemeine Überlegenheit
gegenüber freier Agentenkoordination. Die Realisierung mit MAF
erläutert Abschnitt~\ref{sec:realisierung-agent-executor-workflow}.

Die Abschnitte~\ref{sec:evidenzgebundene-verarbeitung}
bis~\ref{sec:durable-ausfuehrung-agentische-bedienung} vertiefen nun
die sechs Säulen aus Abschnitt~\ref{sec:konsolidierte-designanforderungen}.
Sie erläutern jeweils den Beitrag zum beschriebenen Gesamtablauf.

% Historische Inline-Herkunftsnotizen vor M17; aktuelle Änderung im M17-Bericht.
% [v05-Rest 03.09.: qualitativer Anhangs-Verweis → \ref auf den jetzt
%  angelegten Anhang chap5/anhang/anhang-g1-detail.tex (Autor:
%  einbinden, sonst „??").]
% [Zoom-Prinzip der Blaupause: Karte vollständig, Säulen = Zooms.
%  v03: Einstieg geglättet (Kollegen-P7); „Architekturebenen" statt
%  „Verantwortungsebenen" (Kollegen-P5: State ist keine Verantwortung,
%  Kollision mit G2-„Verantwortungsverteilung" vermieden — weicht
%  bewusst vom Blaupausen-Wort ab, Blaupausen-Nachzug klein).]
% [09.09.: Erst-Seed vor den Struktur-Gates; Beleg:
%  CoreBootstrapStageExecutor + PipelineFullWorkflow. Grafik: eigener
%  Erst-Seed-Pfeil zur Speichernaht; nachgelagerte Strukturvorschläge
%  weiterhin über Prüfung, menschliche Autorisierung und Anwendung.]
% [v05/P1: „eine Abkürzung … existiert nicht" war absoluter als die
%  5.6-Grenze erlaubt (die Naht beweist die Gate-Herkunft nicht
%  rückwirkend) — jetzt als Architekturvertrag des vorgesehenen
%  Pfads formuliert.]
% [K9/K10 + E47-02; heutige Klärung: Speichernaht ohne Trigger-Logik,
%  Fortschreibung zustandsgetrieben (Wiedervorlage!); geschlossenes
%  Decision-Outcome-Vokabular → Karte §1b/⑤.]
% [K11–K13 + Autor-Beschluss 02.09. (H an Doku-Sichten); Kollegen-P3
%  (übernommen): lokales Rendern ungegated (deterministisch,
%  Fingerprint), erst der Publish gated — Beleg GithubDocPublish.cs:
%  „bestehendes github-forward-gate"; deckt sich mit Karte §1b/⑥.
%  K12 GEHÄRTET 02.09.: „kein Doku-Rückweg" = Code-Grenze
%  (GithubDocPublish.cs:12 „KEINE Doc-Ernte", Ernte liest NUR Issues)
%  + Publish-Realität: 14 github_doc_projection-Stempel im Live-Core
%  (sourceRun/contentHash/sha). Merker: → \ref auf 5.7.]
% [v02-Kernkorrektur der Karte: NICHT alles durch den Ledger; K1–K6.
%  Brücken-Satz = Kollegen-Vorschlag 02.09. (gemeinsame Box in der
%  Grafik behauptet keine gemeinsame Implementierung). Kollegen-P2
%  (übernommen): Klärungsantwort als WIEDEREINTRITT ausgewiesen —
%  Karte §0 wörtlich: „Auflösung = Wiedereintritt, kein vierter
%  Außen-Eingang".]
% [Zustandsarten-Tabelle der Karte §4; K16/E47-01. Kollegen-P5
%  (übernommen): „Architekturebenen" + „getrennt geführt".
%  Merker: → \ref 5.8.]
% [v05-Mini-Touch 03.09. (Gesamtdurchgang-P0, todo-v02): auf die
%  Konvergenz-Endform vereinheitlicht (wie 4.9-v07 und 5.8) — die
%  frühere Form „bedienen dieselben Entscheidungspunkte" behauptete
%  implizit identisches Fähigkeitsinventar/Gate-Durchlauf; die
%  damalige Verteidigung stützte sich auf den 4.9-Altwortlaut, der
%  jetzt ebenfalls auf Konvergenz steht (Wort-Parität bleibt).]
% [K14/K15; finale Grafik zeigt den Steward mit Diktat-Kante UND
%  dezentem Steuerungspfeil aufs Verarbeitungsband („Fähigkeit wählen ·
%  Gates vermitteln"). Kollegen-P6 TEILWEISE: Grammatik-Fix +
%  „Schreibweg" übernommen; „bedienen dieselben Entscheidungspunkte"
%  BEHALTEN — Kap. 4.9 sagt wörtlich „beide Zugänge bedienen dieselben
%  Gates und Nähte", und K14 belegt Gate-Antworten über den Steward
%  (Gate-Vermittlung ist real, keine Übertreibung).]
% [Merker EINGELÖST 02.09. (5.1-Vorbau): echtes 4.10-Label ist
%  sec:konsolidierte-designanforderungen (geprüft an chap4.10.tex:51);
%  Tabellen: t:designanforderungen-a1-a10 · t:saeulen-uebergabe.
%  Refs 5.3–5.8 nachziehen (todo.md).]
